\chapter{度量空间.}

正如我们已说过的（见第 146 页），度量空间的概念是 M. 弗雷歇于 1906 年引入的，几年之后又由 F. 豪斯多夫在其《集合论》（Mengenlehre）中加以发展。1920 年以后，它获得了重大的重要性：一方面是随着 S. 巴拿赫及其学派关于赋范空间及其在泛函分析中的应用的基本工作（见第 216 页及以下），另一方面是由于赋值概念在算术和代数几何中唤起的兴趣（在那里，关于一个赋值的完备化特别显示出巨大的成效）。

1920-1930 年这一时期产生了莫斯科学派就度量空间的拓扑性质所开展的一系列研究，这些工作特别致力于求得一个给定的拓扑可度量化的必要充分条件。正是这场思想运动引起了对正规空间概念的兴趣，这一概念是蒂策于 1923 年定义的，但其重要作用只是在乌雷松关于连续数值函数延拓的工作{[}314{]}之后才得到承认。除去实变量的函数这一平凡情形，把定义在一个闭集上的连续数值函数延拓到整个空间的问题，第一次（就平面的情形）是由 H. 勒贝格{[}196 d{]}处理的；在乌雷松的决定性结果之前，它已由 H. 蒂策就度量空间解决{[}307{]}。这个问题向函数取值于任意拓扑空间的情形的推广，近年来在代数拓扑学中获得了相当大的重要性。晚近的工作还揭示出，在这类问题中，正规空间的概念不好驾驭，因为它仍然提供太多的“病态”可能性；大多数时候必须用限制更严的仿紧空间概念去代替它，这一概念是 J. 迪厄多内于 1944 年引入的{[}90 a{]}；在这门理论中，最出色的结果是 A. H. 斯通{[}300{]}证明的那条定理：每个可度量化的空间都是仿紧的。\(^{1}\)

我们已经把注意力引到（见第 155 页）19 世纪末和 20 世纪初的那些重要工作（E. 博雷尔、贝尔、勒贝格、奥斯古德、W. H. 杨）上，这些工作涉及空间 \(R^{n}\) 中点集的分类，以及从连续函数出发、（就函数序列）反复过渡到极限而得到的数值函数的分类和刻画。人们很快看清，度量空间为这类研究提供了一个自然的框架，这类研究在 1910 年之后的发展主要归功于俄国学派和波兰学派。正是这些学派共同揭示了贫集的概念、以及完备度量空间中处处稠密的开集的可数交的那条定理在现代分析中所起的基本作用，这条定理最先（彼此独立地）由奥斯古德{[}240{]}就数值直线、由贝尔{[}11 b{]}就诸空间 \(R^{n}\) 给出证明。

另一方面，1917 年，苏斯林{[}289{]}在纠正勒贝格的一个错误时证明了博雷尔集的连续像不必是博雷尔集，这引导他去定义和研究此后称为“解析集”或“苏斯林集”的那个更大的集类；苏斯林早逝之后，这项研究主要由 N. 卢津（他的思想曾是苏斯林工作的灵感来源）和波兰数学家们（见{[}209{]}和{[}189 b{]}）继续进行。这些集现今的重要性主要系于它们在积分理论中的应用（在那里，多亏它们的特殊性质，一些在任意可测集上不可能实现的构造得以实现），以及现代位势论，在那里，G. 绍凯新近证明的关于苏斯林集的容量的基本定理{[}63{]}，已经显示出它富有各式各样的应用。
% semantic-fragments: legacy-input-seam
\relax{}
